% !TEX program = xelatex
% 공통수학2 · Ⅱ. 집합과 명제 · 중단원(2.명제)·대단원(Ⅱ) 마무리
% 학생용 활동지 (답안 없음)
\documentclass[11pt,a4paper]{article}
\usepackage{kotex}
\usepackage{iftex}
\ifPDFTeX\else
  \setmainhangulfont{Noto Serif CJK KR}
  \setsanshangulfont{Noto Sans CJK KR}
\fi
\usepackage[margin=2.1cm]{geometry}
\usepackage{parskip}
\usepackage{enumitem}
\usepackage{array}
\usepackage{tabularx}
\usepackage{xcolor}
\usepackage{amssymb}
\usepackage{tikz}
\usepackage[most]{tcolorbox}
\usepackage{fancyhdr}
\usepackage{titlesec}

\definecolor{goldc}{HTML}{B07C22}
\definecolor{goldstrong}{HTML}{8A5F16}
\definecolor{indigoc}{HTML}{2B3B57}
\definecolor{indigosoft}{HTML}{6E82A6}
\definecolor{linec}{HTML}{D6DACB}
\definecolor{surface2}{HTML}{E4E8DC}
\definecolor{writeline}{HTML}{C7CCBA}
\definecolor{inksoft}{HTML}{52604F}

\titleformat{\section}{\normalfont\sffamily\bfseries\large\color{indigoc}}{\thesection}{0.6em}{}
\titlespacing{\section}{0pt}{14pt}{6pt}
\newtcolorbox{goalbox}{enhanced, breakable, colback=white, colframe=goldc, boxrule=0.5pt, leftrule=3pt, arc=2pt, left=10pt, right=10pt, top=8pt, bottom=8pt}
\newtcolorbox{sheetbox}{enhanced, breakable, colback=white, colframe=linec, boxrule=0.6pt, arc=3pt, left=11pt, right=11pt, top=9pt, bottom=9pt}
\newcommand{\writespace}[1]{%
  \par\nobreak\vspace{3pt}
  \foreach \n in {1,...,#1} {%
    \noindent\textcolor{writeline}{\rule{\linewidth}{0.4pt}}\par\vspace{0.62cm}%
  }
}
\newcommand{\fillblank}[1]{\makebox[#1]{\hrulefill}}
\pagestyle{fancy}\fancyhf{}\renewcommand{\headrulewidth}{0pt}
\fancyfoot[L]{\footnotesize\color{inksoft} 공통수학2 · 2026학년도 2학기 · 지평선고등학교}
\fancyfoot[R]{\footnotesize\color{inksoft} 다음 시간 --- Ⅲ. 함수와 그래프 시작}

\begin{document}

\noindent
\begin{minipage}[b]{0.62\linewidth}
  {\sffamily\footnotesize\color{indigosoft} 공통수학2 · Ⅱ. 집합과 명제}\\[2pt]
  {\sffamily\bfseries\Large 중단원(2.명제)\textperiodcentered 대단원(Ⅱ) 마무리 \textperiodcentered\ 32차시}
\end{minipage}\hfill
\begin{minipage}[b]{0.36\linewidth}
  \raggedleft\small\color{inksoft}
  반\ \fillblank{1.6cm} \quad 번호\ \fillblank{1.6cm}\\[4pt]
  이름\ \fillblank{3.6cm}
\end{minipage}

\vspace{4pt}
\noindent{\color{goldc}\rule{\linewidth}{2.2pt}}
\vspace{10pt}

\begin{goalbox}
\textbf{오늘의 목표}\\
명제와 조건, 진리집합, 역과 대우, 충분·필요조건을 종합적으로 점검하고, 대단원 Ⅱ `집합과 명제'를 마무리한다.
\vspace{4pt}
\small\color{inksoft}
$\square$ 자신 있음 \quad $\square$ 조금 헷갈림 \quad $\square$ 처음 봄 \hfill (수업 전 나의 상태에 표시)
\end{goalbox}

\section{생각 열기 --- 나만의 개념 지도}

\begin{sheetbox}
`집합과 명제' 대단원에서 배운 개념들을 자유롭게 적고, 서로 관련 있는 것끼리 화살표로 연결해 보자.\\
\small\color{inksoft} (예시 개념: 집합의 연산, 명제와 조건, 진리집합, 역과 대우, 충분·필요조건, 대우증명법, 귀류법, 절대부등식)
\writespace{6}
\end{sheetbox}

\section{종합 문제}

\begin{sheetbox}
\textbf{1.} 다음에서 명제를 모두 찾고, 그 참·거짓을 판별하시오.\\
\hspace*{1.6em}⑴ $\sqrt4$는 유리수이다. \quad ⑵ $x$는 100보다 큰 수이다. \quad ⑶ 26은 3의 배수이다. \quad ⑷ 우리나라는 아름다운 나라이다.
\writespace{3}

\vspace{6pt}
\textbf{2.} 전체집합 $U=\{x\mid x$는 7 이하의 자연수$\}$에 대하여 두 조건 $p$, $q$가\\
\hspace*{1.6em}$p:x^2-6x+8=0$, \quad $q:x^2-6x+9>0$\\
일 때, 다음 조건의 진리집합을 구하시오.\\
\hspace*{1.6em}⑴ $p$ \quad ⑵ $q$ \quad ⑶ $\sim p$ \quad ⑷ $\sim q$
\writespace{4}

\vspace{6pt}
\textbf{3.} 다음 명제의 역과 대우를 말하고, 그 참·거짓을 판별하시오.\\
\hspace*{1.6em}⑴ $x$가 자연수이면 $x$는 유리수이다.\\
\hspace*{1.6em}⑵ $x(x+1)=0$이면 $x+1=0$이다.
\writespace{6}

\vspace{6pt}
\textbf{4.} 두 조건 $p$, $q$가 다음과 같을 때, $p$는 $q$이기 위한 어떤 조건인지 말하시오.\\
\hspace*{1.6em}⑴ $p:x\le1$ \quad $q:-3\le x\le0$\\
\hspace*{1.6em}⑵ $p:a=b$ \quad $q:a^2=b^2$\\
\hspace*{1.6em}⑶ $p:|x|\le1$ \quad $q:-1\le x\le1$
\writespace{5}
\end{sheetbox}

\section{사고의 가시화 --- Connect · Extend · Challenge}

\noindent{\footnotesize\color{inksoft} 대단원 Ⅱ `집합과 명제' 전체를 떠올리며 아래 세 칸을 채워 보자.}\\[6pt]
\begin{tabularx}{\linewidth}{@{}XXX@{}}
\textbf{\footnotesize\color{indigoc} Connect --- 무엇과 연결되나?} &
\textbf{\footnotesize\color{indigoc} Extend --- 어디까지 확장됐나?} &
\textbf{\footnotesize\color{indigoc} Challenge --- 아직 궁금한 것은?} \\[4pt]
\end{tabularx}
\noindent
\begin{minipage}[t]{0.32\linewidth}\writespace{3}\end{minipage}\hfill
\begin{minipage}[t]{0.32\linewidth}\writespace{3}\end{minipage}\hfill
\begin{minipage}[t]{0.32\linewidth}\writespace{3}\end{minipage}

\section{마무리 성찰}

\begin{sheetbox}
\begin{minipage}[t]{0.48\linewidth}
{\footnotesize\color{inksoft} 대단원 Ⅱ `집합과 명제'를 마치며 한 문장으로}
\writespace{2}
\end{minipage}\hfill
\begin{minipage}[t]{0.48\linewidth}
{\footnotesize\color{inksoft} 감사 일기 / 유념 한 줄}
\writespace{2}
\end{minipage}
\end{sheetbox}

\end{document}
